
\section{Introduction}

Hybrid quantum-classical computation is any computation that mixes quantum and classical components.
% 
Many algorithms rely on the principles of mixed computation to work.
% 
For instance, repeat-until-success algorithms involve retrying a state preparation procedure
until a measurement succeeds. 
% 
Specialised languages are required to express hybrid computation, 
as care must be taken to handle the quantum and classical parts properly and separately. 
% 
In the course of my research, I found that languages that address this problem properly are scarce. 
% 
OpenQASM3 \cite{Cross_2022_OpenQASM3} includes some basic quantum feedback mechanisms, but it is difficult to formally study,
as it lacks a formally specified semantics.
% 
Quil \cite{Smith_2016_Quil} does have some formal specification as a low-level quantum PL, but the semantics given is incomplete.
% 
In contrast, Silq \cite{Bichsel_2020_Silq} has a formal semantics and a type system which gives guarantees of safe qubit usage,
as well as mechanisms for working with classical data.
% 
However, no compiler for Silq currently exists, and the complexity of the language makes 
an implementation in the near future unlikely.
% 
The highest standard of formal specification can be found in languages like QWIRE \cite{Paykin_2017_QWIRE, Rand_2019_ReQWIRE},
which is formally verified in Rocq \cite{rocq}.
%
However, QWIRE is not well suited to discussing hybrid programs, as it lacks control primitives like conditions or loops.


Further, I found that a range of possible expressivenesses exist, based on the kinds of control flow one permits in the language.
% 
No semantic work has been done on this stratification, which I believe could be of great use in practise.
% 
Mechanised semantics is a well understood technique, and has been used to build verified compilers
for classical languages (CompCert C \cite{Leroy_2016_compcert}) and even quantum languages (VoQC \cite{Hietala_2023_voqc}).
% 
However, the work in verified quantum compilation is much less well developed than its counterpart in the classical world.
% 
In particular, VoQC works only on circuit optimisation, and cannot deal with more complex or expressive hybrid languages.
% 
I believe this is an important problem to address.


\subsection{Thesis Statement}

\vspace{5mm}

\noindent\fbox{
    \parbox{\textwidth}{
        By giving formal semantics to a family of languages for quantum computers, related by the
        addition of simple features (e.g. individual commands), we will provide safe,
        mathematically grounded implementations of compiler transformations and analyses for
        quantum programs.
    }
}

\subsection{Research goals}

In my work thus far, I have surveyed the state of the art in quantum programming language research, and identified a key gap.
% 
The semantics of low-level languages, particularly hybrid low-level languages with classical control flow, is not well studied. 
% 
I have begun work mechanising semantics for a family of languages, with varying expressive capabilities. 
% 
The aim is to provide a strong foundation for future research, particularly for 
example in verifying optimisations. 
% 
Thus, I derive the following research questions:

\begin{enumerate}[label={\textbf {RQ\arabic*}}]
    \item \label{goal-semantics-simple} Can we extend our work mechanising semantics for quantum programming languages, 
    towards more expressive languages?
    \item \label{goal-optimisation} Can we implement an ``optimisation lifting'' system?
    This should allow one to specify and verify optimisations on more primitive languages in the family, 
    and \emph{lift them up} such that we may get a likewise verified optimisation on the more powerful languages.
    How should the correctness of such a system be specified?
    \item \label{goal-optimisation-specific} How can we leverage this semantic framework to implement verified optimisation for hybrid programs?
    % 
    Can we verify a collection of optimisations (likely as rewrite rules) on the primitive languages 
    and demonstrate their effectiveness on the more powerful members of the language family empirically?
\end{enumerate}

\subsection{Report Structure}
The rest of the report will be structured as follows. 
% 
First, in \cref{background} I will cover a brief, general background in quantum computing.
% 
Moving on, in \cref{work} I will give an overview of my work to date, 
the primary form of this being a lean formalisation of the semantics of a family of quantum programming languages.
% 
I will then move on to discuss my plan for the remainder of my Ph.D. in \cref{plan}, discussing how I will
address the goals I set out above. 
% 
I will briefly touch on other activities tangential 
to my research but highly relevant to my own personal development in \cref{outreach}. 
% 
Finally I will conclude the report in \cref{conclusion}.
